\uuid{G0uO}
\exo7id{3474}
\titre{exo7 3474}
\auteur{quercia}
\organisation{exo7}
\datecreate{2010-03-10}
\isIndication{false}
\isCorrection{true}
\chapitre{Déterminant, système linéaire}
\sousChapitre{Système linéaire, rang}
\cpge{mpsi-mp2i.matrices-applications-lineaires.b, pcsi.matrices-applications-lineaires.b, ptsi.matrices-applications-lineaires.b, mpsi-mp2i.determinants.e, pcsi.determinants.d, mpsi-mp2i.calcul-algebrique-trigonometrie.d, pcsi.calcul-algebrique-trigonometrie.d, ptsi.calcul-algebrique-trigonometrie.d, bcpst1.matrices, bcpst2.applications-lineaires-matrices.b}
\module{Algèbre}
\niveau{L2}
\difficulte{}

\contenu{
\question{
Soit $M = \Bigl(\cos (i+j-1)\theta \Bigr) \in \mathcal{M}_n(\R)$.
Déterminer $\mathrm{rg} M$ en fonction de $\theta$.
}
\reponse{
$M = \Re\left[\begin{pmatrix} e^{i\theta} \cr \vdots \cr e^{ni\theta} \cr \end{pmatrix}
              \Bigl(\begin{matrix} 1 & \dots & e^{(n-1)i\theta} \cr \end{matrix}\Bigr)\right]
  \Rightarrow  \mathrm{rg} M \le 2$.

Le premier mineur $2\times2$ vaut $-\sin^2\theta
 \Rightarrow  \mathrm{rg} M = 2$ si $\theta \not \equiv 0 (\mathrm{mod}\, \pi)$.
Sinon, $\mathrm{rg} M = 1$.
}
}
